\section{Introduction}\label{sec:intro}

For GPU kernels, it is important to have a scheme for describing how to move
tiles of tensors between global memory, shared memory and registers. The de
facto vocabulary for such schemes is the CuTe layout algebra underlying
NVIDIA's CUTLASS library~\cite{cutedocs,cecka}. The most fundamental concept
is the layout, which we shall reuse heavily.

\begin{definition}[Layout]\label{def:layout}
A \emph{shape} is a positive integer $n$, or a tuple $(S_1,\dots,S_k)$ of
shapes; its \emph{size} is $\size n = n$ and
$\size (S_1,\dots,S_k) = \prod_i \size S_i$. A \emph{coordinate} of $S$ is
an integer $i \in [0,n)$, respectively a tuple $(c_1,\dots,c_k)$ of
coordinates of the $S_i$; write $\mathrm{Coord}(S)$ for the set of them. A
\emph{layout} $S{:}D$ is a shape $S$ together with a \emph{stride} $D$ of the
same structure: an integer where $S$ is a positive integer, and a tuple
$(D_1,\dots,D_k)$ where $S$ is $(S_1,\dots,S_k)$. Its index function
$\mathrm{Coord}(S) \to \N$ is defined by recursion on the shape,
\[
(n{:}d)(i) = i\,d,
\qquad
\big((S_1,\dots,S_k){:}(D_1,\dots,D_k)\big)(c_1,\dots,c_k)
= \sum_{j=1}^{k} (S_j{:}D_j)(c_j).
\] Flattening is a bijection $\canon{\cdot}{S} : \mathrm{Coord}(S) \to
[0,\size S)$, taken row-major, so that the last component varies fastest, with
inverse $\mathrm{unflatten}_S$.
\claude{Give the formula here, e.g.\ for $S=(n_1,\dots,n_k)$, $c_k = i \bmod n_k$, $c_{k-1} = \lfloor i/n_k \rfloor \bmod n_{k-1}$, \dots. It is used in Definition~\ref{def:compose} but first written out in Section~\ref{sec:decide}.} A coordinate may therefore be given as a single
integer, and we say that integer \emph{is} the coordinate. The \emph{cosize} of the layout is the least $N$ with its
image inside $[0,N)$, which is one plus the largest value it takes. The layout
is \emph{dense} when its image is an initial segment $[0,N)$ of the integers,
and a \emph{dense bijection} when it is moreover injective, in which case
$N = \size S$ and every value below $N$ is taken exactly once. A
\emph{mode} is one entry of the shape together with the matching entry of the
stride. Two layouts \emph{concatenate}: for $A = S{:}D$ and $B = S'{:}D'$, the
layout $(A,B)$ is $(S,S'){:}(D,D')$, and it sends $(c,c')$ to $A(c) + B(c')$.
\end{definition}

Layouts are fundamentally the correct concept here, because physical memory is
linearly addressed while the logical space is multidimensional and dense as we
imagine it. They are a natural way to abstract and describe those mappings.

The algebra has four operations. Three of them, composition, complement and
logical division, are formalized by Shah~\cite{shah}; the fourth, logical
product, we take from the CUTLASS implementation.
\claude{Notation switches here without warning: Definition~\ref{def:layout} writes $S{:}D$ indexed from 1 and flattens row-major, while the complement formula and the left-divisibility conditions below use Shah's $N{:}d$, indexed from 0, in CuTe's first-fastest order. One sentence saying so, pointing to Section~\ref{sec:cute}, would fix it.}
The \emph{complement} of a
layout
$B = (N_0,\dots,N_\alpha){:}(d_0,\dots,d_\alpha)$, sorted by stride, with
respect to an integer $M$, is
\[
B^{*}_{M} \;=\;
\Big(d_0,\ \tfrac{d_1}{N_0 d_0},\ \dots,\ \tfrac{M}{N_\alpha d_\alpha}\Big)
\;{:}\;
\Big(1,\ N_0 d_0,\ \dots,\ N_\alpha d_\alpha\Big),
\]
chosen so that the concatenation $(B, B^{*}_{M})$ has size $M$ and its layout
function is a bijection of $[0,M)$ onto itself. \emph{Logical division} is
$A \mathbin{/} B = A \circ (B, B^{*}_{\size A})$, which reindexes $A$ so that
one mode runs inside the tile $B$ selects and the other runs across copies of
it. \emph{Logical product} is $A \otimes B = (A, A^{*}_{M} \circ B)$, where
$M = \size A \cdot \cosize B$, the smallest $M$ for which the composition is
defined. Its
first mode is one copy of $A$, and its second walks $B$ through the addresses
$A$ does not occupy, so the result is $\size B$ copies of $A$ placed in the
pattern $B$ names.

These operations are critical for constructing the layouts one usually wants.
There is, however, a lot of complexity underneath them. Layouts are in general
non-injective and non-surjective onto $[0,M)$, and CuTe always wants to write
the composite as a layout again. This causes at least six condition checks. Two are in the complement, whose
shape entries are quotients and so must divide: $N_{i-1}d_{i-1} \mid d_i$ for
every $i$, and $N_\alpha d_\alpha \mid M$. The other four belong to
composition. Write the flattened shape of $A$ as $M_0 \cdots M_\alpha$. It is \emph{left divisible} by $r$
when $r = M_0 \cdots M_{i-1} \cdot c$ for some $i$, where $c$ is smaller than
$M_i$ and divides it, and \emph{weakly left divisible} when the same holds
except that $c$ need not divide $M_i$. For each mode $(N){:}(r)$ of $B$ this is
three checks, since the prefix must divide $r$, the leftover $c$ must divide
$M_i$, and what remains after $r$ must be weakly left divisible by $N$. One
more is required across modes, that the intervals on which they act be pairwise
disjoint~\cite{shah}. None of these conditions is visible in the CuTe API,
yet all of them are visible failures.
\claude{This paragraph lists the six conditions but never says why they are a problem. Open with the framing: CuTe's operations are partial; each is defined only under arithmetic side conditions that the API does not show, that CUTLASS mostly does not check, and that produce wrong layouts silently when they fail. Also, ``visible failures'' is unsupported and, for the complement, false: its failures are silent (the 2026-09-19 draft measures this against CUTLASS 4.8; see the archive).}

We try to resolve those complexities by restricting what a layout can be.

\paragraph{Explicit index spaces.} The input and the output of a layout each
represent one of three things: a logical index into a tile, a thread-value
index, or a physical offset in memory. Therefore a layout should declare which of the three its
input and its output are, and we allow only the maps worth having between them.
Logical and thread-value indices map into one another in either direction.
Logical indices map into physical offsets. Nothing maps out of a physical
offset. By design, a thread reaches memory by composing its map into logical
space with the storage map.

\paragraph{Requirements that follow from the space.} We allow those maps and
not others because the three spaces do not carry the same obligations. Logical
and thread-value indices have to account for their contents losslessly, so a
map into either must be a dense bijection. Composition can only happen when one
layout's image and the other's domain agree completely, which is natural
because a physical offset is forbidden as a domain. A physical offset carries
far fewer requirements, since gaps in memory are common and sharing one address
between threads, to read the same element, is often what is wanted, so a map
into physical need only be alias-free, meaning that distinct coordinates reach
distinct addresses except where the layout marks a broadcast.

\paragraph{Composition and simplification.} The composite of two layouts is not
a layout in general. But we can decide if it can be written as a layout, and
when it is a layout we emit one weighted sum of the coordinates. When it is
not, we emit the composition literally, at the cost of a longer expression and
more work at run time.
